\section{Introduction}

The pointwise convergence problem for the free Schr\"odinger equation
asks for the Sobolev regularity that guarantees convergence to the
initial data almost everywhere.  We establish the critical Sobolev
case in every spatial dimension at least three.  The representative
of the evolution is specified by Gaussian regularization; its
exceptional spatial set is independent of time.

For $n\geq3$, put
\[
 s_n=\frac{n}{2(n+1)}.
\]
We use the Fourier transform
$\widehat f(\xi)=\int_{\R^n}e^{-ix\cdot\xi}f(x)\,\dd x$, extended
from integrable functions to $L^2$, and the inhomogeneous norm
\[
 \norm{f}_{H^s}^2=(2\pi)^{-n}
 \int_{\R^n}(1+|\xi|^2)^s|\widehat f(\xi)|^2\,\dd\xi.
\]
For $a>0$, let
\begin{equation}\label{m:eq-Gaussian}
 U_af(x,t)=(2\pi)^{-n}\int_{\R^n}
 e^{ix\cdot\xi-it|\xi|^2-a|\xi|^2}\widehat f(\xi)\,\dd\xi.
\end{equation}
The integral is absolutely convergent for $f\in L^2$, by
Cauchy--Schwarz.  At a Lebesgue point of $f$, write
\[
 f^*(x)=\lim_{r\downarrow0}\frac1{|B(x,r)|}
                         \int_{B(x,r)}f(y)\,\dd y.
\]

\Needspace{10\baselineskip}
\begin{theorem}[Endpoint convergence]\label{m:main}
For every fixed integer $n\geq3$ and every complex-valued
$f\in H^{s_n}(\R^n)$, there is a full-measure set $E_f$ of Lebesgue
points of $f$ and a function $v(x,t)$, $x\in E_f$, $0\leq t<1$, such
that $v(x,\cdot)$ is continuous, $v(x,0)=f^*(x)$, and
\begin{equation}\label{m:eq-uniform-Gaussian}
 \lim_{a\downarrow0}\sup_{0<t<1}|U_af(x,t)-v(x,t)|=0
 \qquad(x\in E_f).
\end{equation}
In particular the finite Gaussian limit
\[
 u_f(x,t)=\lim_{a\downarrow0}U_af(x,t)
\]
exists for every real $0<t<1$ on this one set, and
$\lim_{t\downarrow0}u_f(x,t)=f^*(x)$.
\end{theorem}

The Gaussian limit is taken first.  No Fourier support, radiality,
logarithmic strengthening of $H^{s_n}$, or additional regularity is
assumed.  The theorem concerns all real times approaching zero, not a
chosen sequence.

\subsection{History and significance}

Carleson proved the one-dimensional convergence theorem for
$H^{1/4}(\R)$, and Dahlberg and Kenig proved failure below that
regularity \citep{Carleson1980,DahlbergKenig1982}.  The
one-dimensional threshold therefore includes its endpoint.
Sj\"olin and Vega independently obtained convergence for $s>1/2$
in higher dimensions \citep{Sjolin1987,Vega1988}.  At $s>1/2$,
Sj\"ogren and Sj\"olin also constructed representatives continuous
on whole time lines outside one exceptional spatial set, using spatial
or spacetime averages \citep[Theorem~2]{SjogrenSjolin1989}.  In the plane,
restriction and bilinear advances by Bourgain, Moyua--Vargas--Vega,
Tao--Vargas and Tao
\citep{Bourgain1995,MoyuaVargasVega1996,TaoVargas2000,Tao2003}
preceded Lee's $s>3/8$ theorem \citep{Lee2006}.  In dimensions
$n\geq3$, Bourgain obtained the sufficient range
$s>1/2-1/(4n)$ \citep{Bourgain2013}.

Bourgain's later counterexample gives the necessary condition
$s\geq n/(2(n+1))$ in dimensions $n\geq2$
\citep[Proposition~1]{Bourgain2016}.  That obstruction does not
assert failure at equality.  Du, Guth and Li proved convergence
for $s>1/3$ in dimension two
\citep[Theorems~1.1--1.2]{DuGuthLi2017}.  Du, Guth, Li and Zhang
then proved higher-dimensional convergence for
$s>(n+1)/(2(n+2))$ and multilinear refined Strichartz estimates
\citep[Theorems~1.1 and~4.2]{DuGuthLiZhang2018}.  Du and Zhang
proved convergence for $s>n/(2(n+1))$ in higher dimensions
\citep[Theorem~1.1]{DuZhang2019}.  Recent work improves
higher-dimensional maximal bounds in other $L^p$ ranges
\citep[Theorems~1.1--1.2]{DuLi2024}; a Bessel-potential
extension still assumes $s>n/(2(n+1))$ at $p=2$
\citep[Theorems~1.3--1.4]{PanSunTan2026}.

\Cref{m:main} resolves the critical Sobolev case of Carleson's
pointwise convergence problem in higher dimensions.  The distinction
between equality and a strict Sobolev inequality is decisive: an
arbitrarily small uncompensated frequency loss does not prove the
endpoint theorem.  We use the planar companion,
\emph{Endpoint convergence for the planar Schr\"odinger equation}
\citep[Sections~3 and~9]{PSE2026}, specifically its
\emph{Limiting information calculus} and \emph{Block comparison and
telescoping costs} lemmas and its time-tree mechanism.
Its planar conclusion is not invoked as a higher-dimensional result.  The
dimension-dependent frame, projection and curvature arguments needed
here are established below, with the hypotheses of the reusable
lemmas stated at their points of use.

The analytic background includes positive-curvature decoupling
\citep{BD2015}, multilinear refined Strichartz estimates
\citep[Theorem~4.2]{DuGuthLiZhang2018}, and the refined-packet line of
Guth--Iosevich--Ou--Wang and Du--Iosevich--Ou--Wang--Zhang
\citep{GIOW2020,DIOWZ2021}.  Singh and Singh Parmar identified a
support-versus-counting collar mismatch in the normalized printed
arbitrary-packet formulation and proved a buffered correction in
ambient dimension two
\citep[Theorems~1.1--1.2 and Appendix~A]{SinghSinghParmar2026}.
\Cref{p:buffered-refined} establishes the needed buffered
activity-count estimate in every ambient dimension used below,
following that induction structure and using decoupling; neither the
printed arbitrary-packet
estimate nor the planar correction is imported as an all-dimensional
input.  Other inputs are multilinear Kakeya and restriction
\citep{BCT2006,Guth2010,Tao2020}, discretized projections
\citep{He2020}, and fractal Schr\"odinger estimates
\citep{DuZhang2019}.  The multilinear restriction estimate is used
away from its endpoint, and decoupling losses are removed through
strict induction gains; an endpoint estimate without such losses is
not assumed.

\Needspace{9\baselineskip}
\subsection{The new estimates and the proof}

There are three principal steps.

\paragraph{A transverse frame estimate.}
In $k$ spatial dimensions, \Cref{f:frame} counts unit cubes where
$k+1$ transverse packet fields are simultaneously large relative to
their square envelopes, allowing a separate witness for each field.
Its constants depend polynomially on inverse transversality and on
the fixed geometric parameters.  The estimate applies to compact
packet profiles and is therefore stable under the later geometric
reductions.  Its proof combines weighted packet inequalities with a
finite hierarchy of amplitudes and occupancies.  A positive smoothing
argument handles gaps that are subpower in the original scale, while
a robust projection theorem forces the limiting amplitude-ratio profiles.
The explicit restricted-path construction keeps the probability laws
compatible with cancellations and direction-dependent deletions.

\paragraph{A geometric induction in full dimension.}
The fractal estimates and the full-dimensional transverse saving first
yield a sparse gain.  An induction on
packet scale and sparsity then reduces a large level set to a
polynomial wall.  Codimension-two angular narrowing alone does not
exclude a curved ruled wall.  We also prove a quadratic-sublevel
narrow estimate and use it to detect the second fundamental form.
Coefficient-uniform graph estimates and paths inside the small-Hessian
set convert the relevant wall pieces into affine plates.  Flattening
along these plates makes \Cref{f:frame} applicable in dimension $n-1$.
Two energy budgets, one per time bin and one global with density
weights, ensure that summing the geometric pieces loses no growing
factor.

\paragraph{A common time tree.}
The induction produces a weak $L^p$ maximal bound for some $p>2$ on
time intervals of length $1/N$ at frequency $N$, with precisely the
factor $N^{s_n}$.  Write $S(t)f=e^{it\Delta}f$ for the $L^2$ evolution
with multiplier $e^{-it|\xi|^2}$.  Using one measurable time graph for all frequencies,
the time-tree argument gives frequency square summation and the local
bound
\begin{equation}\label{m:eq-roadmap-maximal}
 \int_\Omega\sup_{0\leq t<1}|S(t)f(x)|\,\dd x
 \lesssim\norm{f}_{H^{s_n}},
\end{equation}
initially for compact Fourier support, uniformly on translated unit
balls $\Omega$.  Uniform-in-time approximation then produces a common
representative.  Applying the Gaussian approximate identity to
countably many error envelopes proves
\eqref{m:eq-uniform-Gaussian} without an uncountable intersection of
exceptional sets.  This also proves \Cref{m:main}; the details appear
in \Cref{e:trace}.

The packet estimates lead to the frame theorem, \Cref{f:frame};
separately, the full-dimensional fractal and transverse estimates,
including \Cref{t:transverse}, lead to the sparse estimate,
\Cref{s:sparse}.  These two inputs meet in \Cref{w:weak}: the wall
geometry and flattening developed there make the frame theorem
applicable with $k=n-1$.  The final section proves the time-tree and
Gaussian representative assertions.  All constants may depend on the fixed
dimension.  Whenever a sufficiently large finite profile order is
required, it is chosen once for the relevant array class; no induction
increases that order.
